\section{Clase 10 - Pr\'actica 5 (M\'etodos iterativos para sistemas lineales)}

Si las matrices son muy grandes no conviene usar m\'etodos directos. Estudiaremos los m\'etodos de Jacobi y Gauss - Seidel.

$$ 
	\begin{cases}
		4x_{1} + x_{2} = 5\\
		2x_{1} + 3x_{2} = 5\\
    \end{cases}
    \rightarrow \
    \begin{cases}
		4x_{1} = 5 - x_{2}\\
		3x_{2} = 5 - 2x_{1}\\
    \end{cases}
$$

Empiezo con cierto $x^{(0)} \in \mathbb{R}^{2}$ y propongo:

$$ 
    \begin{cases}
		4x_{1}^{(k + 1)} = 5 - x_{2}^{(k)}\\
		3x_{2}^{(k + 1)} = 5 - 2x_{1}^{(k)}\\
    \end{cases}
    \ \text{es decir:}\ 
    \begin{cases}
		x_{1}^{(k + 1)} = \frac{1}{4}(5 - x_{2}^{(k)})\\
		x_{2}^{(k + 1)} = \frac{1}{3}(5 - 2x_{1}^{(k)})\\
    \end{cases}
    \ \text{Jacobi}
$$

$$ 
    \begin{cases}
		4x_{1}^{(k + 1)} = 5 - x_{2}^{(k)}\\
		3x_{2}^{(k + 1)} = 5 - 2x_{1}^{(k + 1)}\\
    \end{cases}
    \ \text{es decir:}\ 
    \begin{cases}
		x_{1}^{(k + 1)} = \frac{1}{4}(5 - x_{2}^{(k)})\\
		x_{2}^{(k + 1)} = \frac{1}{3}(5 - 2x_{1}^{(k + 1)})\\
    \end{cases}
    \ \text{Gauss - Seidel}
$$

Con Jacobi: $
	\begin{pmatrix}
        0\\
        0\\
	\end{pmatrix}
	\rightarrow
	\begin{pmatrix}
        \frac{5}{4}\\
        \frac{5}{3}\\
	\end{pmatrix}
	\rightarrow
	\begin{pmatrix}
        \frac{10}{12}\\
        \frac{5}{6}\\
	\end{pmatrix}
	\cdots
$\\

Con Gauss - Seidel: $
	\begin{pmatrix}
        0\\
        0\\
	\end{pmatrix}
	\rightarrow
	\begin{pmatrix}
        \frac{5}{4}\\
        \frac{5}{6}\\
	\end{pmatrix}
	\rightarrow
	\begin{pmatrix}
        \frac{25}{24}\\
        \frac{35}{36}\\
	\end{pmatrix}
	\cdots
$\\

Qu\'e pasa si los elementos de la diagonal son nulos? Si $A$ es inversible, siempre puedo reordenar ecuaciones e inc\'ognitas para que queden coeficientes no nulos en la diagonal.\\

\textit{Ejemplo}:

$$
	\begin{pmatrix}
        0 & 1 & 2\\
        2 & 0 & 3\\
        0 & 1 & 0\\
	\end{pmatrix}
	\rightarrow_{
	F_{1} \leftrightarrow F_{2}}
    \begin{pmatrix}
        2 & 0 & 3\\
        0 & 1 & 2\\
        0 & 1 & 0\\
	\end{pmatrix}
	\rightarrow_{
	C_{2} \leftrightarrow C_{3}}
	\begin{pmatrix}
        2 & 3 & 0\\
        0 & 2 & 1\\
        0 & 0 & 1\\
	\end{pmatrix}
$$

$$ 
	\begin{cases}
		x_{1} + 2x_{3} = b_{1}\\
		2x_{1} + 3x_{3} = b_{2}\\
		x_{2} = b_{3}\\
    \end{cases}
    \
    \begin{cases}
		2x_{1} + 3x_{3} = b_{2}\\
		x_{1} + 2x_{3} = b_{1}\\
		x_{2} = b_{3}\\
    \end{cases}
    \
    \begin{cases}
		2\tilde{x}_{1} + 3\tilde{x}_{2} = b_{2}\\
		\tilde{x}_{3} + 2\tilde{x}_{2} = b_{1}\\
		\tilde{x}_{3} = b_{3}\\
    \end{cases}
$$

\underline{En general}:

\subsection{M\'etodo de Jacobi}

$$ 
	\begin{cases}
		x^{0} \in \mathbb{R}^{n}\\
		x_{i}^{(k + 1)} = \frac{1}{a_{ii}}\bigg(b_{i} - \sum \limits_{j = i,\ j \neq i}^{n} a_{ij}x_{j}^{(k)} \bigg),\ 1 \leq i \leq n\\
    \end{cases}
$$

\subsection{M\'etodo de Gauss - Seidel}

$$ 
	\begin{cases}
		x^{0} \in \mathbb{R}^{n}\\
		x_{i}^{(k + 1)} = \frac{1}{a_{ii}}\bigg(b_{i} - \sum \limits_{j = i}^{i - 1} a_{ij}x_{j}^{(k + 1)} - \sum \limits_{j = i + 1}^{n} a_{ij}x_{j}^{(k)} \bigg),\ 1 \leq i \leq n\\
    \end{cases}
$$\\

Si escribimos $A = L + D + U$, con $L$ y $U$ triangulares inferior y superior estrictas de $A$, respectivamente y $D$ diagonal de $A$, entonces:

$$Ax = b \iff Dx = -(L + U)x + b \iff (L + D)x = -Ux + b$$

\begin{itemize}
	\item \underline{Jacobi}:
	
	$$Dx^{(k + 1)} = -(L + U)x^{(k)} + b$$
	
	$$\Rightarrow x^{(k + 1)} = -D^{-1}(L + U)x^{(k)} + D^{-1}b$$
	
	\item \underline{Gauss - Seidel}:
	
	$$(L + D)x^{(k + 1)} = -Ux^{(k)} + b$$
	
	$$\Rightarrow x^{(k + 1)} = -(L + D)^{-1}Ux^{(k)} + (L + D)^{-1}b$$
\end{itemize}

\subsection{M\'etodos iterativos en general}

Proponemos:

$$x^{(k + 1)} = Bx^{(k)} + C$$

\begin{itemize}
	\item Jacobi: $B_{J} = -D^{-1}(L + U),\ C_{J} = D^{-1}b$
	
	\item Gauss - Seidel: $B_{GS} = -(L + D)^{-1}U,\ C_{GS} = (L + D)^{-1}b$
\end{itemize}

\underline{Definici\'on}: Decimos que el m\'etodo iterativo converge si $\ \forall \ x^{0} \in \mathbb{R}^{n}\ \exists \ x\ / \ \lim \limits_{k \rightarrow \infty} x^{(k)} = x$.\\

Si $B$ y $C$ son tales que resolver $Ax = b$ es equivalente a resolver $x = Bx + C$ tenemos que si el m\'etodo es convergente, converge a una soluci\'on de $Ax = b$:

$$x^{(k + 1)} = Bx^{(k)} + C \rightarrow x = Bx + C$$

\underline{Error}: Sea $x^{*}\ / \ x^{*} = Bx^{*} + C$:

$$e^{(k + 1)} = x^{(k + 1)} - x^{*} = Bx^{(k)} + C - (Bx^{*} + C) = B(x^{(k)} - x^{*}) = Be^{(k)} = \cdots = B^{k + 1}e^{(0)}$$

Luego,

$$\|e^{(k + 1)}\|' = \|B^{k + 1}e^{(0)}\|' \leq \|B^{k + 1}\|'\|e^{(0)}\|' \leq \|B\|'^{k + 1}\|e^{(0)}\|'$$

\underline{Teorema}: Sea $x^{*}$ tal que $x^{*} = Bx^{*} + C$. Si $\|B\|' < 1$ para alguna norma matricial subordinada a una norma vectorial, la iteraci\'on de punto fijo dada por $x^{(k + 1)} = Bx^{k} + C$ converge a $x^{*}$. Es decir, $\lim \limits_{k \rightarrow \infty} x^{(k)} = x^{*}\ \forall \ x^{0}$ inicial.\\

En el ejemplo inicial:

$$ A =
	\begin{pmatrix}
        4 & 1\\
        2 & 3\\
	\end{pmatrix}
	=
    \begin{pmatrix}
        0 & 0\\
        2 & 0\\
	\end{pmatrix}
	+
	\begin{pmatrix}
        4 & 0\\
        0 & 3\\
	\end{pmatrix}
	+
	\begin{pmatrix}
        0 & 1\\
        0 & 0\\
	\end{pmatrix}
	= L + D + U
$$

$$ B_{J} = -D^{-1}(L + U) = -
	\begin{pmatrix}
        \frac{1}{4} & 0\\
        0 & \frac{1}{3}\\
	\end{pmatrix}
    \begin{pmatrix}
        0 & 1\\
        2 & 0\\
    \end{pmatrix}
	=
	\begin{pmatrix}
        0 & -\frac{1}{4}\\
        -\frac{4}{3} & 0\\
	\end{pmatrix}
$$

$$ B_{GS} = -(D + L)^{-1}U = -
	\begin{pmatrix}
        \frac{1}{4} & 0\\
        -\frac{1}{6} & \frac{1}{3}\\
	\end{pmatrix}
    \begin{pmatrix}
        0 & 1\\
        0 & 0\\
    \end{pmatrix}
	=
	\begin{pmatrix}
        0 & -\frac{1}{4}\\
        0 & \frac{1}{6}\\
	\end{pmatrix}
$$

$$\Rightarrow \|B_{J}\|_{\infty} = \frac{2}{3} < 1,\ \|B_{GS}\|_{\infty} = \frac{1}{4} < 1$$

y ambos m\'etodos convergen.\\

\underline{Cuidado}: Consideramos $B =
	\begin{pmatrix}
        \frac{1}{2} & 10^{6}\\
        0 & \frac{1}{2}\\
	\end{pmatrix}
$. Vale que $\|B\|_{\infty} = \|B\|_{1} = 10^{6} + \frac{1}{2}$. Sin embargo:

$$ B = \frac{1}{2}
	\begin{pmatrix}
        1 & 2\cdot10^{6}\\
        0 & 1\\
	\end{pmatrix}
	,\ B^{2} = \frac{1}{4}
    \begin{pmatrix}
        1 & 4\cdot10^{6}\\
        0 & 1\\
	\end{pmatrix}
	, \cdots,\ B^{k} = \frac{1}{2^{k}}
	\begin{pmatrix}
        1 & 2k\cdot10^{6}\\
        0 & 1\\
	\end{pmatrix}
$$

y por lo tanto, $B^{k} \rightarrow_{k \rightarrow \infty} 
	\begin{pmatrix}
        0 & 0\\
        0 & 0\\
	\end{pmatrix}
$\\

Si $B$ es real y sim\'etrica, entonces $B$ es diagonal. Existe una matriz inversible $\tilde{C}$ tal que:

$$\tilde{C}^{-1}B\tilde{C} =
	\begin{pmatrix}
        \lambda_{1} & \cdots & \cdots\\
        \vdots & \vdots & \vdots\\
        \cdots & \cdots & \lambda_{n}\\
	\end{pmatrix}
	= \tilde{D}
$$

Luego,

$$B^{k} = \tilde{C}\tilde{D}^{k}\tilde{C}^{-1} = \tilde{C}
	\begin{pmatrix}
        \lambda_{1}^{k} & \cdots & \cdots\\
        \vdots & \vdots & \vdots\\
        \cdots & \cdots & \lambda_{n}^{k}\\
	\end{pmatrix}
	\tilde{C}^{-1}
$$

y vale que:

$$B^{k} \rightarrow_{k \rightarrow \infty} 0 \iff |\lambda_{j}|^{k} \rightarrow_{k \rightarrow \infty} 0 \iff |\lambda_{j}| < 1\ \forall \ j \iff \rho(B) < 1$$

\underline{Teorema}: Sea $A \in \mathbb{R}^{n \times n}$, entonces $\rho(A) = \inf \limits_{\|.\|'} \|A\|'$, donde el \'infimo se toma sobre todas las normas matriciales subordinadas a normas vectoriales.\\

\underline{Teorema de Schur}: Sea $A \in \mathbb{C}^{n \times n}$, existe una matriz unitaria $Q$ ($Q^{-1} = \bar{Q}^{t} = Q^{*}$) y una matriz triangular superior $T$ con los autovalores de $A$ en la diagonal tal que $Q^{*}AQ = T$.\\

\underline{Demostraci\'on}: Si $n = 1$, es obvio. Supongamos que el resultado vale para matrices de tama\~no $(n - 1) \times (n - 1)$. Sea $u$ un autovector de $A$ de norma 1 ($\|u\|_{2} = 1$) asociado al autovalor $\lambda$.\\

Consideremos $\{u, w_{2}, \cdots, w_{n}\}$ B.O.N. de $\mathbb{C}^{n}$, y sea $\tilde{Q}$ cuyas columnas son dicha base ortogonal. Entonces:

$$\tilde{Q}^{*}A\tilde{Q} = \tilde{Q}^{*}A\bigg(u\bigg|w_{2}\bigg| \cdots \bigg|w_{n}\bigg) = \tilde{Q}^{*}\bigg(\lambda u\bigg|Aw_{2}\bigg| \cdots \bigg|Aw_{n}\bigg) = \bigg(\lambda \tilde{Q}^{*}u\bigg|\tilde{Q}^{*}Aw_{2}\bigg| \cdots \bigg|\tilde{Q}^{*}Aw_{n}\bigg)$$

$$=
	\begin{pmatrix}
        \lambda_{1} & z_{1} & \cdots & z_{n - 1}\\
        0 &  &  & \\
        \vdots &  & A^{(1)} &  \\
        0  &  &  & \\
	\end{pmatrix}
$$

con $A^{(1)} \in \mathbb{R}^{(n - 1) \times (n - 1)}$ y $\tilde{Q}^{*} = \begin{pmatrix}
        u^{1} & \cdots & \cdots & u^{n}\\
        w_{2}^{1} & \cdots & \cdots & w_{2}^{n}\\
        \vdots & \vdots & \vdots & \vdots\\
        w_{n}^{1} & \cdots & \cdots & w_{n}^{n}\\
	\end{pmatrix}
$. Luego, existen $Q^{(1)}$ unitaria y $T^{(1)}$ triangular superior con los autovalores de $A^{(1)}$ en la diagonal tal que $(Q^{(1)})^{*}A^{(1)}Q^{(1)} = T^{(1)}$. Sea $M =
	\begin{pmatrix}
        1 & 0 & \cdots & 0\\
        0 &  &  & \\
        \vdots &  & Q^{(1)} &  \\
        0  &  &  & \\
	\end{pmatrix}
$ y sea $Q = \tilde{Q}M$:

$$Q^{*}AQ = M^{*}\tilde{Q}^{*}A\tilde{Q}M = 
	\begin{pmatrix}
        1 & 0 & \cdots & 0\\
        0 &  &  & \\
        \vdots &  & (Q^{(1)})^{*} &  \\
        0  &  &  & \\
	\end{pmatrix}
	\begin{pmatrix}
        \lambda_{1} & z_{1} & \cdots & z_{n - 1}\\
        0 &  &  & \\
        \vdots &  & A^{(1)} &  \\
        0  &  &  & \\
	\end{pmatrix}
	\begin{pmatrix}
        1 & 0 & \cdots & 0\\
        0 &  &  & \\
        \vdots &  & Q^{(1)} &  \\
        0  &  &  & \\
	\end{pmatrix}
$$

$$\Rightarrow Q^{*}AQ = 
	\begin{pmatrix}
        \lambda_{1} & z_{1} & \cdots & z_{n - 1}\\
        0 &  &  & \\
        \vdots &  & (Q^{(1)})^{*}A^{(1)} &  \\
        0  &  &  & \\
	\end{pmatrix}
	\begin{pmatrix}
        1 & 0 & \cdots & 0\\
        0 &  &  & \\
        \vdots &  & Q^{(1)} &  \\
        0  &  &  & \\
	\end{pmatrix}
	=
	\begin{pmatrix}
        \lambda_{1} &  & z^{t}Q^{(1)} & \\
        0 &  &  & \\
        \vdots &  & T^{(1)} &  \\
        0  &  &  & \\
	\end{pmatrix}
	= T
$$

Como los autovalores de $T$ son los autovalores de $A$ y son los elementos de su diagonal, tenemos el resultado buscado.\\

\underline{Idea de la demostraci\'on del teorema anterior}\\

Queremos ver que

\begin{itemize}
	\item $\rho(A) \leq \|A\|'$ para toda norma matricial $\|.\|'$ subordinada a una norma vectorial.
	
	\item dado $\epsilon > 0$, existe una norma $\|.\|'_{\epsilon}$ vectorial tal que la norma matricial subordinada cumple $\|A\|_{\epsilon} \leq \rho(A) + \epsilon$.
\end{itemize}

Veamos esto \'ultimo: Por Schur, existe una matriz $Q$ unitaria tal que $Q^{*}AQ = T$, donde $T$ es triangular superior con los autovalores de $A$ en la diagonal. Sea $D = 
	\begin{pmatrix}
        \delta & \cdots & \cdots & 0\\
        0 & \delta^{2} & \cdots & 0\\
        \vdots & \vdots & \vdots & \vdots\\
        0 & \cdots & 0 & \delta^{n}\\
	\end{pmatrix}
$, entonces:

$$D^{-1}TD = 
	\begin{pmatrix}
        \frac{1}{\delta} & \cdots & \cdots & 0\\
        0 & \frac{1}{\delta^{2}} & \cdots & 0\\
        \vdots & \vdots & \vdots & \vdots\\
        0 & \cdots & 0 & \frac{1}{\delta^{n}}\\
	\end{pmatrix}
	TD
	=
	\begin{pmatrix}
        \frac{\lambda_{1}}{\delta} & \frac{t_{21}}{\delta} & \cdots & \frac{t_{1n}}{\delta}\\
        0 & \frac{\lambda_{2}}{\delta^{2}} & \cdots & 0\\
        \vdots & \vdots & \vdots & \vdots\\
        0 & \cdots & 0 & \frac{\lambda_{n}}{\delta^{n}}\\
	\end{pmatrix}
	\begin{pmatrix}
        \delta & \cdots & \cdots & 0\\
        0 & \delta^{2} & \cdots & 0\\
        \vdots & \vdots & \vdots & \vdots\\
        0 & \cdots & 0 & \delta^{n}\\
	\end{pmatrix}
$$

$$
	=
	\begin{pmatrix}
        \lambda_{1} & t_{12}\delta & \cdots & t_{1n}\delta^{n - 1}\\
        0 & \lambda_{2} & \cdots & t_{2n}\delta^{n - 2}\\
        \vdots & \vdots & \vdots & \vdots\\
        0 & \cdots & 0 & \lambda_{n}\\
	\end{pmatrix}
	= \tilde{D} + \tilde{T}
	,\ \text{con}\ T\ \text{estrictamente triangular superior.}
$$

Sea $S = QD$, luego:

$$S^{-1}AS = D^{-1}Q^{*}AQD = D^{-1}TD = \tilde{D} + \tilde{T}$$

$$\|\tilde{T}\|_{\infty} = \max \limits_{1 \leq i \leq n} \sum \limits_{j = 1}^{n} |\tilde{t}_{ij}| = \max \limits_{1 \leq i \leq n} \sum \limits_{j > i}^{n} |\tilde{t}_{ij}\delta^{j - i}| \leq \bigg( \max \limits_{1 \leq i \leq n} \sum \limits_{j > i}^{n} |\tilde{t}_{ij}| \bigg)\delta \leq \|T\|_{\infty}\delta,\ 0 < \delta < 1$$

$$\Rightarrow \|S^{-1}AS\|_{\infty} = \|\tilde{D} + \tilde{T}\|_{\infty} \leq \|\tilde{D}\|_{\infty} + \|\tilde{T}\|_{\infty} = \max \limits_{1 \leq j \leq n} |\lambda_{j}| + \|\tilde{T}\|_{\infty} \leq \rho(A) + \|T\|_{\infty}\delta$$

Si tomamos $\delta < \frac{\epsilon}{\|T\|_{\infty}}$, tenemos $\|S^{-1}AS\|_{\infty} < \rho(A) + \epsilon$.\\

Definimos la norma vectorial $\|.\|_{\epsilon}$ en $\mathbb{C}^{n}$ como $\|v\|_{\epsilon} = \|S^{-1}v\|_{\infty}$ (verificar que es una norma en $\mathbb{C}^{n}$ y que induce una norma en $\mathbb{R}^{n}$). Se obtiene la correspondiente norma matricial subordinada. Entonces:

$$\|A\|_{\epsilon} = \max \limits_{x \in \mathbb{R}^{n} - \{0\}} \frac{\|Ax\|_{\epsilon}}{\|x\|_{\epsilon}} = \max \limits_{x \in \mathbb{R}^{n} - \{0\}} \frac{\|S^{-1}Ax\|_{\infty}}{\|S^{-1}x\|_{\infty}} = \max \limits_{x \in \mathbb{R}^{n} - \{0\}} \frac{\|S^{-1}ASS^{-1}x\|_{\infty}}{\|S^{-1}x\|_{\infty}} \leq \max \limits_{y \in \mathbb{C}^{n} - \{0\}} \frac{\|S^{-1}ASy\|_{\infty}}{\|y\|_{\infty}}$$

$$= \|S^{-1}AS\|_{\infty} \leq \rho(A) + \epsilon$$

como quer\'iamos probar.